\chapter{Git}
\label{ch:git}
\chapterlead{Git 记录文件的历史，也是协作的接口。本章从第一次提交写到提交一个补丁，重点是三件事：理解三个区域、写出可读的提交信息、在出错时能安全地退回来。}

\section{起步：从零到第一次提交}

\begin{shellsession}
$ mkdir -p ~/lab/repo && cd ~/lab/repo
$ git init -b main
$ git config user.name "Weiwen Zhang"
$ git config user.email "weiwen@example.com"
$ echo "# 实验仓库" > README.md
$ git add README.md
$ git commit -m "初始化仓库"
$ git log --oneline
\end{shellsession}

配置分三级：\texttt{-{}-global} 写入 \texttt{\textasciitilde/.gitconfig}，\texttt{-{}-system} 影响全机，不加参数只对当前仓库生效。在 NixOS 上，用户级配置由 Home Manager 生成更稳：

\begin{nixcode}
programs.git = {
  enable = true;
  userName = "Weiwen Zhang";
  userEmail = "weiwen@example.com";
  extraConfig = {
    init.defaultBranch = "main";
    pull.rebase = true;
    core.pager = "delta";
  };
};
\end{nixcode}

\section{三个区域}

\begin{figure}[H]
\centering
\figfit{%
\begin{tikzpicture}[font=\small, node distance=7mm,
  box/.style={draw=ink!45, fill=ink!5, rounded corners=1pt, inner sep=5pt, align=center},
  arr/.style={-{Latex[length=1.5mm]}, teal, thick}]
  \node[box] (work) {工作区\\[-1pt]{\scriptsize 正在编辑的文件}};
  \node[box, right=of work] (index) {暂存区\\[-1pt]{\scriptsize 下次提交的内容}};
  \node[box, right=of index] (repo) {仓库\\[-1pt]{\scriptsize 提交历史}};
  \draw[arr] ([yshift=2mm]work.east) -- node[above, font=\scriptsize] {git add} ([yshift=2mm]index.west);
  \draw[arr] ([yshift=2mm]index.east) -- node[above, font=\scriptsize] {git commit} ([yshift=2mm]repo.west);
\end{tikzpicture}%
}
\caption{三个区域。提交前进暂存区，是把“这次改了什么”明确下来的动作}
\end{figure}

理解暂存区的意义在于选择：一次改了五个文件，但只想提交其中两个，\texttt{git add} 指定的那两个即可。这让提交保持主题单一。

\begin{shellsession}
$ git status            # 哪些文件变了，处在哪个区域
$ git diff              # 工作区与暂存区的差异
$ git diff --staged     # 暂存区与上次提交的差异
$ git add -p file       # 逐块选择要暂存的部分
\end{shellsession}

\section{提交信息}

一条可读的提交信息由两行组成：标题与空行后的正文。

\begin{textcode}
修复备份脚本在目标目录不存在时的失败

原脚本假设目标目录已存在，首次运行会因 rsync 报错退出。
改为先创建目录，并在失败时打印目标路径。
\end{textcode}

标题写做了什么，不超过一行；正文写为什么，以及改动的边界。标题不要写“更新文件”“修改代码”这类无法检索的句子。半年后回看历史时，能靠标题判断是否与当前问题相关。

\section{分支与合并}

\begin{shellsession}
$ git switch -c fix/backup
$ git switch main
$ git merge fix/backup
$ git branch -d fix/backup
\end{shellsession}

\texttt{switch -c} 新建并切换分支。分支只是一条指向提交的名字，创建成本近乎为零。合并有两种：\texttt{merge} 保留分支结构，\texttt{rebase} 把提交重放到目标分支之上，历史成为一条直线。

个人项目用 rebase 保持线性，团队协作中已推送的分支不 rebase，因为它会改写别人已经拉取的提交。

\section{撤销的四种情况}

\begin{booktable}{\textwidth}{J{0.889} J{1.111}}
\toprule
情况 & 做法 \\
\midrule
\endhead
改了文件但没 add & \texttt{git restore 文件} \\
add 了但没 commit & \texttt{git restore -{}-staged 文件} \\
提交了但没推送 & \texttt{git commit -{}-amend} 或 \texttt{git reset -{}-soft HEAD\textasciitilde1} \\
已经推送了 & \texttt{git revert 提交}，生成一个反向提交 \\
\bottomrule
\end{booktable}

已经推送的历史不要用 \texttt{reset} 改写，因为它会让其他人的仓库与远端分叉。\texttt{revert} 是安全的做法：新增一个提交抵消旧提交的效果，历史完整保留。

\begin{warnbox}[危险命令]
\texttt{git reset -{}-hard} 会丢弃工作区改动，\texttt{git clean -fd} 会删除未跟踪文件，两者都不可恢复。执行前先运行 \texttt{git status} 确认将失去什么，或先 \texttt{git stash} 暂存。
\end{warnbox}

\section{提交一个补丁}

给开源项目提交改动的流程基本固定。

\begin{shellsession}
$ git clone https://github.com/owner/project.git
$ cd project
$ git switch -c fix/my-change
$ git add -p
$ git commit
$ git format-patch -1
$ git push fork fix/my-change
\end{shellsession}

\texttt{format-patch} 生成可直接用邮件发送的补丁文件，许多项目仍以此接收贡献。用平台提交时改为推送分支并发起合并请求。

提交前应当跑一遍项目自己的检查：格式化、静态检查、测试。这些命令通常写在项目的 \texttt{README} 或 \texttt{CONTRIBUTING} 里。补丁因为格式问题被退回，是最容易避免的一类返工。

\section{练习}

\begin{enumerate}
  \item 建立仓库，提交三个主题互不相同的改动，用 \texttt{git log --oneline} 检查标题是否可读。
  \item 修改三个文件，只暂存其中两个，提交后用 \texttt{git show -{}-stat} 确认范围。
  \item 用 \texttt{git restore} 撤销工作区改动，用 \texttt{git restore -{}-staged} 撤销暂存，记录两步的差别。
  \item 制造一次提交错误，用 \texttt{-{}-amend} 修正标题，再用 \texttt{git reflog} 找到修正前的提交。
  \item 为一个真实项目写一份补丁，跑通项目自带的检查。
\end{enumerate}

\begin{checkpoint}
本章结束时，应能：
\begin{itemize}
  \item 说清工作区、暂存区、仓库三者的关系；
  \item 写出含标题与正文的提交信息，并说明正文写什么；
  \item 用分支隔离改动，判断何时用 merge、何时用 rebase；
  \item 按四种情况选择正确的撤销方式；
  \item 完成从分支到补丁的完整贡献流程。
\end{itemize}
\end{checkpoint}
